%قبل از شروع مسئله ترجیحاً چیزی ننویسید. اما اگر خواستید بنویسید، در قالب لیستی (فقرات و شمارش) بنویسید.
%برای اضافه کردن نمره، بعدی از شروع‌مسئله، [x] را اضافه کنید و به جای x نمره مورد نظر را قرار دهید. برای زیرمسئله‌ها هم به همین صورت است.
%\شروع‌مسئله[10]
\شروع‌مسئله [12]
%\متن‌سیاه{سوال عملی اول:}\\
%\متن‌سیاه{دراپ‌تست}\\

%\متن‌سیاه{ورودی:}\\
%در خط اول ورودی، ابتدا k و سپس m داده می‌شود.\\
%\متن‌سیاه{خروجی:}\\
%در تنها سطر خروجی، تعداد آیفون‌های مورد نیاز می‌آید. اگر در هیچ صورتی چنین کاری ممکن نبود ۱- چاپ می‌شود.\\
%\متن‌سیاه{محدودیت‌ها:}\\
%$1\leq k\leq 10^4$

%$m\leq10^9$
%\\
%\textbf{\متن‌سیاه{ورودی و خروجی نمونه}\\
%% تست اول
%\begin{latin}
% \begin{tabularx}{\textwidth}{|X | X|}
% \hline
% \textbf{stdin} &
% \textbf{stdout}\\
% \hline
% 100 100&
% 1\\
% \hline
% 
% \end{tabularx}
%\end{latin}
%
%% تست دوم
%
%\begin{latin}
% \begin{tabularx}{\textwidth}{|X | X|}
% \hline
% \textbf{stdin} &
% \textbf{stdout}\\
% \hline
% 340 20&
% 3\\
% \hline
% 
% \end{tabularx}
%\end{latin}}
% صورت اصلی مسئله را در این‌جا بنویسید.
در هر یک از موارد زیر تعدادی الگوریتم داریم که زمان اجرایشان بر اساس دو ورودی $n$ و $k$ داده شده است. برای هر حالت $n$ و $k$ تعیین کنید بهتر است از کدام الگوریتم استفاده شود (شرط‌ها تا جای ممکن ساده باشند).

%اگر مسئله چند قسمت داشت به این صورت عمل کنید:
\شروع{زیرمسائل}
\زیرمسئله[3] % مسئلۀ ۱آ
الگوریتم اول:
$\Theta(nk^{2})$
و 
الگوریتم دوم:
$\Theta(n^{2}k)$
\زیرمسئله[3] % مسئلۀ ۱ب
الگوریتم اول:
$\Theta(n^{4}k^{6})$
و 
الگوریتم دوم:
$\Theta(n^{5}k^{3})$
\زیرمسئله[3] % مسئلۀ ۱پ
الگوریتم اول:
$\Theta(\frac{n^{3}k^{3}}{\log{n}})$
و 
الگوریتم دوم:
$\Theta(n^{5}k\log^{2}{n})$
\زیرمسئله[3] % مسئلۀ ۱ت
الگوریتم اول:
$\Theta(n^{3}\sqrt{k})$
، 
الگوریتم دوم:
$\Theta(n^{2}k^{2})$
و 
الگوریتم سوم:
$\Theta(k^{4})$
\پایان{زیرمسائل}

\solution{
شروط استفاده هر کدام از الگوریتم‌ها به ترتیب در زیر نوشته شده:
%اگر مسئله چند قسمت داشت به این صورت عمل کنید:
\شروع{فقرات}
\فقره مورد اول
الگوریتم اول:
$k \leq n$
و 
الگوریتم دوم:
$n \leq k$

\فقره مورد دوم
الگوریتم اول:
$k^{3} \leq n$
و 
الگوریتم دوم:
$n \leq k^{3}$

\فقره مورد سوم
الگوریتم اول:
$k \leq n\log^{\frac{3}{2}}n$
و 
الگوریتم دوم:
بقیه حالات

\فقره مورد چهارم
الگوریتم اول:
$n \leq k^{\frac{7}{6}}$
، 
الگوریتم دوم:
در هیچ حالتی الگوریتم بهینه نیست
و 
الگوریتم سوم:
$k^{\frac{7}{6}} \leq n$
\پایان{فقرات}

}




%%%%%%%%%%%%%%%%%%%%%%%%%%%%%%%%%%%%%%%%%%%%%%%%%%%%%
% برخی دستورات مفید برای تمرین‌ها. %
%%%%%%%%%%%%%%%%%%%%%%%%%%%%%%%%%%%%%%%%%%%%%%%%%%%%%

% فرمول‌های نمونه:
%$f(x) = \Theta(x)$
%$T(x, y) \leq \log(x) + 2^y + \binom{2n}{n}$
% {\tt code\_function}


% لیست موارد بدون شماره‌گذاری:
%\شروع{فقرات}
% \فقره مورد اول 
% \شروع{فقرات}
% \فقره مورد اول داخلی 
% \شروع{فقرات}
% \فقره اولین مورد خیلی داخلی 
% \فقره دومین مورد خیلی داخلی 
% \پایان{فقرات}
% \فقره مورد دوم داخلی 
% \پایان{فقرات}
% \فقره مورد دوم 
%\پایان{فقرات}

% لیست موارد شماره‌گذاری شده:
%\شروع{شمارش}
% \فقره یک مورد 
% \فقره یک مورد دیگر 
% \فقره باز هم یک مورد دیگر
%\پایان{شمارش}

% برای نوشتن یک فرمول چند خطی که یک شماره بخوره. برای حالت‌های دیگر به سایت زیر مراجعه کنید:
% https://www.overleaf.com/learn/latex/Aligning_equations_with_amsmath
%\begin{align} 
% \begin{split}
% (x+y)^3 &= (x+y)^2(x+y) \\
% &= (x^2+2xy+y^2)(x+y) \\
% &= (x^3+2x^2y+xy^2) + (x^2y+2xy^2+y^3) \\
% &= x^3+3x^2y+3xy^2+y^3
% \end{split} 
%\end{align}

% نوشتن ماتریس و آرایه دوبعدی
%\begin{align}
% A = 
% \begin{bmatrix}
% A_{11} & A_{21} \\
% A_{21} & A_{22}
% \end{bmatrix}
%\end{align}

% اضافه کردن تصویر با شرح. برای اضافه کردن تصویر بدون شرح، فقط دستور درج‌تصویر را استفاده کنید.
%\شروع{شکل}[h!]
%\تنظیم‌ازوسط
% \درج‌تصویر[پهنا=0.5\پهنای‌متن]{img.jpg}
%\شرح{شرح تصویر}
%\برچسب{تصویر:۱}
%\پایان{شکل}
%\پایان{زیرمسائل}

%کد زیر را ببینید.
%%\lstinputlisting{source_filename.py}
%\begin{latin}
%\begin{lstlisting}
%# Import the modules
%import sys
%import random
%
%ans = True
%
%while ans:
%question = raw_input("Ask the magic 8 ball a question: (press enter to quit) ")
%
%answers = random.randint(1,8)
%
%if question == "":
%sys.exit()
%
%elif answers == 1:
%print "It is certain"
%
%elif answers == 2:
%print "Outlook good"
%
%elif answers == 3:
%print "You may rely on it"
%
%elif answers == 4:
%print "Ask again later"
%
%elif answers == 5:
%print "Concentrate and ask again"
%
%elif answers == 6:
%print "Reply hazy, try again"
%
%elif answers == 7:
%print "My reply is no"
%
%elif answers == 8:
%print "My sources say no"
%\end{lstlisting}
%\end{latin}